\documentclass[11pt]{article}
\usepackage[a4paper,margin=25mm]{geometry}
\usepackage{amsmath,amssymb,amsthm,microtype}
\usepackage[colorlinks=true,urlcolor=blue,citecolor=blue]{hyperref}
\newtheorem{theorem}{Theorem}
\newtheorem{lemma}[theorem]{Lemma}
\newtheorem{corollary}[theorem]{Corollary}
\newcommand{\E}{\mathbb E}
\newcommand{\Pp}{\mathbb P}
\newcommand{\Var}{\operatorname{Var}}
\newcommand{\Law}{\mathcal L}
\title{Polynomially long consecutive failures of log-concavity\\in the bulk of tree independence sequences}
\author{Research draft prepared for Jan Mikul\'ik with ChatGPT}
\date{27 September 2026}
\begin{document}
\maketitle
\begin{abstract}
For each integer $b\ge64$, we give a fixed family of trees of maximum
degree $b+1$ whose independence sequences contain a consecutive run of
$\Omega(N^{\log_b(bq)})$ strict log-concavity failures at a limiting
relative position strictly below $1$. Here $q<1/2$ is an explicitly
specified algebraic number. At maximum degree $65$, the run has length
$\Omega(N^{4/5})$ and lies below $0.99$ times the independence number.
Allowing the fixed degree bound to depend on $\varepsilon$ yields
$\Omega(N^{1-\varepsilon})$ consecutive bulk failures for every
$\varepsilon>0$. Consequently the extremal run-length exponent is $1$.
The analytic step is a finite-height, stretched-exponential Fourier bound;
it gives uniform convergence of an exact lattice interpolation and all its
derivatives. This closes the transfer from a limiting mass valley to a
consecutive block of finite coefficient inequalities.
\end{abstract}

\noindent\textbf{Status.} This is a proposed proof with all analytic steps
supplied below, not an independently refereed result or a formal
verification. No claim of bibliographic priority is made. The lower-bound
constants and the starting height in this stronger run theorem are not
numerically evaluated.

\section{Statement}
For a tree $T$, let $i_k(T)$ count its independent sets of size $k$, and
write $\alpha(T)$ for its independence number. A failure is an interior
index $k$ such that $i_k(T)^2<i_{k-1}(T)i_{k+1}(T)$.

Fix an integer $b\ge64$. Let $T_{b,0}$ be an edge rooted at one endpoint.
Construct $T_{b,h+1}$ by attaching a new root to the roots of $b$ disjoint
copies of $T_{b,h}$. Put
\begin{equation}\label{eq:size}
 N_h=b^h+\frac{b^{h+1}-1}{b-1},\qquad
 \alpha_h=b^h+\sum_{j=0}^{\lfloor(h-1)/2\rfloor}b^{h-1-2j}.
\end{equation}
The sum is empty at $h=0$. These are respectively the order and
independence number of $T_{b,h}$; its maximum degree is at most $b+1$.

Let $\delta$ be the unique solution in $(2^{-b},2^{1-b})$ of
\begin{equation}\label{eq:parameters}
 \delta(2-\delta)^b=1,
 \qquad r=1-\delta,\quad q=\frac{r}{1+r},\quad \rho=bq,
 \quad \pi=\frac{q}{1+q}.
\end{equation}
Define
\begin{equation}\label{eq:beta}
 A_b=1+\frac{b}{b^2-1},\qquad
 L_b=r+\delta\pi+\frac{\pi}{b-1},\qquad
 \beta_b=\frac{L_b}{A_b},\qquad \kappa_b=\log_b\rho.
\end{equation}

\begin{theorem}\label{thm:main}
For every fixed integer $b\ge64$, there exist constants $c_b>0$ and
$H_b<\infty$ such that, for every $h\ge H_b$, there is an interval of
consecutive interior indices $I_h$ with
\begin{align}
 |I_h|&\ge c_b\rho^h=\Omega_b(N_h^{\kappa_b}),\label{eq:run}\\
 i_k(T_{b,h})^2&<i_{k-1}(T_{b,h})i_{k+1}(T_{b,h})
 \quad(k\in I_h),\label{eq:failure}\\
 \sup_{k\in I_h}\left|\frac{k}{\alpha_h}-\beta_b\right|&\longrightarrow0.
 \label{eq:location}
\end{align}
Moreover
\begin{equation}\label{eq:gap}
 0<\beta_b<\frac{(3b-2)(b+1)}{3(b^2+b-1)}
 =1-\frac{2b-1}{3(b^2+b-1)}<1.
\end{equation}
For $b=64$, the entire run lies below $0.99\alpha_h$ for all sufficiently
large $h$, and has length $\Omega(N_h^{4/5})$.
\end{theorem}

\begin{corollary}\label{cor:epsilon}
For every $\varepsilon>0$, there are constants $\Delta_\varepsilon<\infty$
and $d_\varepsilon>0$ and a family of arbitrarily large trees of maximum
degree at most $\Delta_\varepsilon$ with $\Omega_\varepsilon(N^{1-\varepsilon})$
consecutive failures, all at indices below $(1-d_\varepsilon)\alpha(T)$.
The degree bound and relative gap are fixed along each family.
\end{corollary}

Let $M(N)$ denote the largest length of a consecutive run of failures
among trees on \emph{at most} $N$ vertices.
\begin{corollary}\label{cor:exponent}
The extremal logarithmic exponent is exactly one:
\[
 \lim_{N\to\infty}\frac{\log M(N)}{\log N}=1,
 \qquad\text{equivalently }M(N)=N^{1-o(1)}.
\]
This does not assert $M(N)=\Theta(N)$.
\end{corollary}

\paragraph{Context.}
Galvin's Question 3.1 asks for log-concavity failure at a fixed relative
distance from the last coefficient~\cite{Galvin}. Consecutive failures
have been studied in~\cite{Ramos,Yu}. Yu's construction gives
$\Omega(N^{1/3})$ consecutive failures with indices $k/\alpha(T)\to1$;
its discussion explicitly raises a bulk analogue. Fang et al. prove
central log-concavity for sufficiently large forests~\cite{Fang}; the
location here is outside their guaranteed interval. The present argument
strengthens the earlier single-failure construction on the same rooted
trees. It neither changes that construction nor contradicts unimodality.

\section{Exact recursion and its normalized limit}
At $\delta=2^{-b}$ the left side of~\eqref{eq:parameters} is below $1$;
at $\delta=2^{1-b}$ it is at least $2(1-b2^{-b})>1$. Its derivative has
the sign of $2-(b+1)\delta$, which is positive between these endpoints.
Thus the specified solution exists uniquely. In particular
\begin{equation}\label{eq:bounds}
 0<\delta<b^{-2},\quad \frac{49}{100}\le q<\frac12,
 \quad \sqrt b<\rho<b.
\end{equation}
Use the same activity $\lambda=r/\delta=r(1+r)^b$ at every vertex.
An independent set has probability proportional to $\lambda^{|I|}$.
The occupied/unoccupied root partition-function ratio is
$\lambda/(1+\lambda)=r$ for $T_{b,0}$, and the recursion
$R_{h+1}=\lambda/(1+R_h)^b$ makes it exactly $r$ at every height.
Thus each isolated rooted subtree has root occupation probability $q$.

Let $X_h$ be the size under this probability law, let
$\mu_h^s=\E(X_h\mid\text{root}=s)$, and let $m_h=\E X_h$.
Root conditioning gives
\[
 \mu_{h+1}^0=b[(1-q)\mu_h^0+q\mu_h^1],\qquad
 \mu_{h+1}^1=1+b\mu_h^0.
\]
Therefore, with $D_h=\mu_h^1-\mu_h^0$,
\begin{equation}\label{eq:mode}
 D_0=\delta,\quad D_{h+1}=1-\rho D_h,\quad
 D_h=\theta+a_h,\qquad
 \theta=\frac1{1+\rho},\quad a_h=(\delta-\theta)(-\rho)^h.
\end{equation}
Here $\delta<\theta$: indeed $\delta<b^{-2}<1/(1+b/2)<\theta$.
The signed normalization $a_h$ removes the alternating growing mode
exactly. Define the centered conditional variables and a deterministic
correction by
\[
 Y_h^s=\frac{X_h-\mu_h^s}{a_h}\ \Big|\ (\text{root}=s),\qquad
 z_h=\frac{D_h}{a_h}=1+\frac{\theta}{a_h}.
\]
If $S_i$ are independent Bernoulli$(q)$ variables, with conditionally
independent copies of the appropriate conditional laws, then
\begin{equation}\label{eq:smoothing}
 \begin{split}
 Y_{h+1}^0&\ \overset d=\ -\rho^{-1}
   \sum_{i=1}^b\bigl(Y_{h,i}^{S_i}+(S_i-q)z_h\bigr),\\
 Y_{h+1}^1&\ \overset d=\ -\rho^{-1}\sum_{i=1}^bY_{h,i}^0.
 \end{split}
\end{equation}

\begin{lemma}\label{lem:contraction}
The pair $(\Law(Y_h^0),\Law(Y_h^1))$ converges in the maximum of the two
$2$-Wasserstein distances to a unique centered finite-variance fixed pair
$(\Law(Y^0),\Law(Y^1))$ of~\eqref{eq:smoothing} with $z_h$ replaced by $1$.
Both limiting variances $v_s$ are positive and at most $5/(4b)$.
\end{lemma}
\begin{proof}
On pairs of centered finite-second-moment laws, call the transformation
in~\eqref{eq:smoothing} $\mathcal T_z$. Couple two input pairs optimally
in each state, use common $S_i$, and take independent couplings over $i$.
The differences have mean zero, so their sum's second moment is the sum
of their second moments. Consequently
\[
 d(\mathcal T_z\nu,\mathcal T_z\widetilde\nu)
 \le \frac{\sqrt b}{\rho}\,d(\nu,\widetilde\nu),\qquad
 d(\mathcal T_z\nu,\mathcal T_1\nu)
 \le\frac{\sqrt{bq(1-q)}}{\rho}|z-1|.
\]
The space is complete (centering is closed under $W_2$ convergence),
and $\ell:=\sqrt b/\rho<1$. The contraction theorem gives the fixed pair.
Its distance $e_h$ from the height-$h$ pair obeys
$e_{h+1}\le\ell e_h+C\rho^{-h}$, so $e_h\to0$ (indeed
$e_h=O_b(\ell^h)$).

Put $t=b/\rho^2=1/(bq^2)<1$. Taking variances in the fixed equations gives
\[
 v_1=tv_0,\qquad
 v_0=t[(1-q)v_0+qv_1+q(1-q)],
 \qquad
 v_0=\frac{tq(1-q)}{1-t(1-q)-qt^2}>0.
\]
Since $q^2\ge(49/100)^2>6/25$, we have $t<25/(6b)$.
For $b\ge64$,
\[
 1-t(1-q)-qt^2\ge1-t-t^2
 \ge1-\frac{25}{384}-\left(\frac{25}{384}\right)^2>\frac56.
\]
Thus $0<v_1\le v_0\le (25/(24b))/(5/6)=5/(4b)$.
\end{proof}

The unconditional centered normalization is
\begin{equation}\label{eq:Z}
 Z_h:=\frac{X_h-m_h}{a_h}\ \overset d=\ Y_h^S+(S-q)z_h
 \ \longrightarrow\ Z:=Y^S+(S-q)
\end{equation}
in $W_2$, where $S$ has probability $q$ of being $1$.

\section{A Fourier bound uniform at finite height}
Weak convergence alone does not decide neighboring coefficient signs.
We now control the complete fundamental frequency interval before taking
any derivative or lattice difference.

Let $\chi_h^s(u)=\E(e^{iuX_h}\mid\text{root}=s)$ and
$Q_h(u)=\max_s|\chi_h^s(u)|$. The exact conditional recursion is
\begin{equation}\label{eq:fourierrec}
 \chi_{h+1}^0(u)=((1-q)\chi_h^0(u)+q\chi_h^1(u))^b,
 \qquad \chi_{h+1}^1(u)=e^{iu}(\chi_h^0(u))^b.
\end{equation}
In particular, for $1\le k\le h$,
\begin{equation}\label{eq:amplify}
 Q_h(u)\le Q_k(u)^{b^{h-k}}.
\end{equation}

\begin{lemma}\label{lem:fourier}
Set $\gamma=\log b/\log\rho>1$. There are positive constants $K,c$ and
an integer $h_0$, depending only on $b$, such that for every $h\ge h_0$
and $|t|\le\pi|a_h|$,
\begin{equation}\label{eq:domination}
 \max_s|\E e^{itY_h^s}|\le K e^{-c|t|^\gamma},
 \qquad |\E e^{itZ_h}|\le K e^{-c|t|^\gamma}.
\end{equation}
\end{lemma}
\begin{proof}
Let $\psi_s$ be the characteristic functions of $Y^s$.
Positive finite variance gives
$|\psi_s(t)|^2=1-v_st^2+o(t^2)$ at zero. Choose $\eta>0$ sufficiently
small that both moduli are strictly below $1$ throughout
$\eta\le|t|\le\rho\eta$. Convergence in $W_2$ gives uniform convergence
of characteristic functions on compact intervals, since
$|\psi_h^s(t)-\psi_s(t)|\le|t|W_2(\Law(Y_h^s),\Law(Y^s))$.
Thus, after increasing $h_0$ if necessary, there is $u_*<1$ such that
\begin{equation}\label{eq:annulus}
 Q_k(t/a_k)\le u_*\quad
 (k\ge h_0,\ \eta\le|t|\le\rho\eta).
\end{equation}
Increase $h_0$ again to ensure $\eta/|a_{h_0}|<\pi$.

First suppose
$\eta/|a_h|\le|u|\le\eta/|a_{h_0}|$.
Choose $k\in[h_0,h]$ with
$\eta\le|a_ku|\le\rho\eta$; successive values grow by the factor $\rho$.
Writing $t=a_hu$, we have
$|t|\le\rho\eta\rho^{h-k}$. Hence~\eqref{eq:amplify}--\eqref{eq:annulus}
give
\[
 Q_h(u)\le u_*^{b^{h-k}}
 \le\exp\left(-\frac{-\log u_*}{(\rho\eta)^\gamma}|t|^\gamma\right).
\]

For the remaining physical frequencies
$\eta/|a_{h_0}|\le|u|\le\pi$, observe that
\[
 X_1\mid0\ \overset d=\operatorname{Bin}(b,2q),\qquad
 X_1\mid1\ \overset d=1+\operatorname{Bin}(b,r).
\]
Indeed $X_0$ is Bernoulli$(2q)$, while $X_0\mid0$ is Bernoulli$(r)$.
Since $0<2q,r<1$, both conditional laws have lattice span one, and
$Q_1(u)<1$ for $0<|u|\le\pi$. Compactness gives a common bound $v_*<1$
on the remaining interval. Let $a_0=\delta-\theta$. Then
\[
 Q_h(u)\le v_*^{b^{h-1}}
 \le\exp\left(-\frac{-\log v_*}{b(\pi|a_0|)^\gamma}|t|^\gamma\right),
\]
because $|t|\le\pi|a_0|\rho^h$ and $\rho^\gamma=b$.
For $|t|<\eta$ a constant prefactor suffices. Centering does not change
modulus; nor does the sign of $a_h$ affect these bounds. The unconditional
characteristic function is a convex combination of the centered
conditional ones with unit-modulus phase factors. This proves both claims.
\end{proof}

\section{Smooth densities and exact lattice interpolation}
\begin{lemma}\label{lem:densities}
The laws of $Y^0,Y^1,Z$ have densities $f_0,f_1,f$ that extend to entire
functions and are strictly positive everywhere on the real line.
If $\phi_h(t)=\E e^{itZ_h}$ and
\begin{equation}\label{eq:interp}
 g_h(x)=\frac1{2\pi}\int_{-\pi|a_h|}^{\pi|a_h|}
                 e^{-itx}\phi_h(t)\,dt,
\end{equation}
then, for every fixed integer $j\ge0$,
\begin{equation}\label{eq:Cj}
 \sup_{x\in\mathbb R}|g_h^{(j)}(x)-f^{(j)}(x)|\longrightarrow0.
\end{equation}
At each lattice node $x_{h,k}=(k-m_h)/a_h$ one has exactly
\begin{equation}\label{eq:lattice}
 g_h(x_{h,k})=|a_h|\Pp(X_h=k).
\end{equation}
\end{lemma}
\begin{proof}
Lemma~\ref{lem:fourier} passes to the limiting characteristic functions.
Since $\gamma>1$, the bound $Ke^{-c|t|^\gamma}$ remains integrable after
multiplication by $|t|^j e^{R|t|}$ for every finite $j,R$. Fourier
inversion therefore gives densities with entire extensions. On the real
axis they are nonnegative. Neither conditional density is identically
zero, so its real zeros are isolated and it is positive almost everywhere.

The fixed equation for $Y^1$ expresses its density as the scaled and
reflected $b$-fold convolution of $f_0$. A convolution of two integrable
densities positive almost everywhere is strictly positive at every point:
its integrand is positive almost everywhere. These densities are bounded,
so the convolution is finite. Iteration shows $f_1>0$ everywhere.
The density of $Y^S+(S-q)$ is
\[
 f(x)=(1-q)f_0(x+q)+qf_1(x-(1-q))>0.
\]
The fixed equation for $Y^0$ now expresses $f_0$ as a scaled reflected
$b$-fold convolution of this positive density, so $f_0>0$ everywhere as
well. The identities first hold almost everywhere and then everywhere by
continuity.

Extend $\phi_h(t)$ by zero outside its fundamental interval in
\eqref{eq:interp}. Pointwise convergence and Lemma~\ref{lem:fourier}
give, by dominated convergence,
\[
 \int_{\mathbb R}|t|^j
 \left|\phi_h(t)\boldsymbol1_{\{|t|\le\pi|a_h|\}}-\E e^{itZ}\right|dt
 \longrightarrow0.
\]
Fourier inversion and differentiation under the integral prove
\eqref{eq:Cj}. The functions $g_h$ are real by conjugate symmetry.
Finally, substitute $t=a_hu$ in~\eqref{eq:interp} at $x_{h,k}$.
Lattice Fourier inversion over $[-\pi,\pi]$ gives~\eqref{eq:lattice};
the absolute Jacobian $|a_h|$ handles either sign of $a_h$.
\end{proof}

\section{A valley forces a whole interval of positive log curvature}
\begin{lemma}\label{lem:valley}
There are a closed interval $J$ of positive length and a constant
$\zeta>0$ such that $(\log f)''\ge2\zeta$ throughout $J$.
One can take $J$ inside $[-q-1/4,1-q+1/4]$.
\end{lemma}
\begin{proof}
In the mixture $Z=Y^S+(S-q)$, the component means are $-q$ and $1-q$,
their variances are at most $5/(4b)\le5/256$, and both weights are at
least $49/100$. Consider intervals of radius $1/4$ centered at $-q$,
$1/2-q$, and $1-q$, denoted $J_L,J_M,J_R$ in order. Chebyshev's
inequality gives
\[
 \min\{\Pp(Z\in J_L),\Pp(Z\in J_R)\}
 \ge\frac{49}{100}\left(1-\frac5{16}\right)=\frac{539}{1600},
 \qquad
 \Pp(Z\in J_M)\le\frac5{16}=\frac{500}{1600}.
\]
For the middle interval, each component is at least $1/4$ from its own
mean. Endpoints have probability zero.

If a positive continuous density is log-concave on the hull of three
ordered intervals of equal length, the middle mass is at least the
smaller outer mass. To see this, choose a point in each outer interval
whose density is at least that interval's average. Concavity of the
logarithm bounds the density at every intervening point below by the
smaller chosen density, and integration over the middle interval proves
the assertion. The strict inequalities above contradict this property.
Since $f>0$ is smooth, $\log f$ is twice continuously differentiable.
If its second derivative were everywhere nonpositive in the hull's
interior it would be concave there. It therefore has a point of strictly
positive second derivative. Continuity supplies $J$ and $\zeta$.
\end{proof}

By~\eqref{eq:Cj} with $j=0,1,2$, positivity of $f$ on $J$ implies that,
for all sufficiently large $h$, $g_h>0$ and
\begin{equation}\label{eq:curvature}
 (\log g_h)''\ge\zeta\quad\text{on }J.
\end{equation}
Let $I_h$ consist of all integers $k$ whose three nodes
$x_{h,k-1},x_{h,k},x_{h,k+1}$ belong to $J$. These integers are
consecutive, even when $a_h<0$, and
\begin{equation}\label{eq:count}
 |I_h|\ge |J||a_h|-4\ge c_b\rho^h
\end{equation}
for all sufficiently large $h$ and a constant $c_b>0$.

Write $p_{h,k}=\Pp(X_h=k)$ and $s_h=1/|a_h|$. The exact integral
formula for a centered second difference, together with
\eqref{eq:lattice}--\eqref{eq:curvature}, yields
\begin{align}
 \log\frac{p_{h,k-1}p_{h,k+1}}{p_{h,k}^2}
 &=\int_{-s_h}^{s_h}(s_h-|u|)(\log g_h)''(x_{h,k}+u)\,du\notag\\
 &\ge\frac{\zeta}{|a_h|^2}>0\qquad(k\in I_h).\label{eq:finitecurvature}
\end{align}
All three probabilities are positive, by~\eqref{eq:lattice} and
positivity of $g_h$ on $J$. Thus these are interior indices of the true
coefficient support. Finally
$p_{h,k}=i_k(T_{b,h})\lambda^k/Z_{T_{b,h}}(\lambda)$, so
\begin{equation}\label{eq:transfer}
 \frac{i_{k-1}(T_{b,h})i_{k+1}(T_{b,h})}{i_k(T_{b,h})^2}
 =\frac{p_{h,k-1}p_{h,k+1}}{p_{h,k}^2}
 \ge\exp\left(\frac{\zeta}{|a_h|^2}\right)>1.
\end{equation}
This is a finite-height conclusion valid uniformly at every index in the
run. No estimate of the form ``pointwise error is smaller than mesh
squared'' is assumed; the exact interpolation and its $C^2$ convergence
are what justify the signs.

\section{Location and extremal consequences}
We finish the proof of Theorem~\ref{thm:main}. At depth $j$ in the
$b$-ary skeleton the occupation probability $p_j$ satisfies
$p_0=q$, $p_{j+1}=q(1-p_j)$, whence
\[
 p_j=\pi+(q-\pi)(-q)^j.
\]
A terminal skeleton vertex and its extra leaf together have expectation
$r+\delta p_h$. Summing over levels,
\[
 \frac{m_h}{b^h}=r+\delta p_h+\sum_{j=0}^{h-1}b^{j-h}p_j
 \longrightarrow r+\delta\pi+\frac{\pi}{b-1}=L_b.
\]
Formula~\eqref{eq:size} follows alternatively from conditional maxima
$u_0=v_0=1$, $u_{h+1}=b\max(u_h,v_h)$, $v_{h+1}=1+bu_h$;
induction gives the displayed alternating-level sum. Hence
$\alpha_h/b^h\to A_b$. Since $J$ is bounded and
$|a_h|/b^h=|a_0|q^h\to0$, the nodes defining $I_h$ satisfy
\[
 \sup_{k\in I_h}|k-m_h|=O_b(\rho^h),\qquad
 \sup_{k\in I_h}\left|k/\alpha_h-\beta_b\right|\to0.
\]
Since $r+\delta\pi<1$ and $\pi<1/3$, we have
\[
 \beta_b<\frac{1+1/[3(b-1)]}{1+b/(b^2-1)}
 =\frac{(3b-2)(b+1)}{3(b^2+b-1)}<1.
\]
For $b=64$ this upper bound is $12350/12477<99/100$.
Also $\rho\ge64(49/100)>30$ and $30^5>64^4$, so
$\kappa_{64}>4/5$. This proves every assertion of the theorem.

As $b\to\infty$, $q\to1/2$ and
$\kappa_b=1+\log q/\log b\to1$. Given $\varepsilon>0$, fix an integer
$b\ge64$ with $\kappa_b>1-\varepsilon$ and use
$d_\varepsilon=(1-\beta_b)/2$ and $\Delta_\varepsilon=b+1$.
This proves Corollary~\ref{cor:epsilon}. Here $b$ is fixed before
$h\to\infty$; no estimate uniform in $b$ is being asserted.

For Corollary~\ref{cor:exponent}, fix the same $b$. The orders satisfy
$N_{h+1}=bN_h+1$. For every sufficiently large budget $N$, choose the
largest $h$ with $N_h\le N$. Then $N_h>(N-1)/b$, and therefore
\[
 M(N)\ge c'_bN_h^{\kappa_b}
 \ge c''_bN^{\kappa_b}.
\]
It follows that $\liminf\log M(N)/\log N\ge1-\varepsilon$ for every
$\varepsilon>0$. The trivial upper bound $M(N)\le N$ completes the
proof.\hfill$\square$

\paragraph{What is and is not quantified.}
The tree family, activity, exponent and limiting bulk location are
explicit. The argument establishes positive run constants and a finite
starting height for each fixed $b$, and proves the required uniform
passage to finite coefficients. It does not evaluate the curvature
interval $J$, its lower curvature bound, or the resulting $H_b$.
In particular this paper does not assert that a polynomially long run
already occurs at the height $8$ used in the earlier single-failure
certificate. The inequalities here may be very weak in relative size:
the guaranteed logarithmic ratio in~\eqref{eq:finitecurvature} is of
order $\rho^{-2h}$.

\begin{thebibliography}{9}
\bibitem{Galvin} D. Galvin, \emph{Trees with non log-concave independent
set sequences}, arXiv:2502.10654v2 (2026), Question 3.1.
\url{https://arxiv.org/abs/2502.10654}.
\bibitem{Ramos} C. Bautista-Ramos, C. Guill\'en-Galv\'an and
A. G\'omez-Salgado, \emph{Linear Recurrences for Non-Log-Concave Independence
Polynomials of Trees}, Graphs and Combinatorics 42, 59 (2026).
\url{https://doi.org/10.1007/s00373-026-03054-4}.
\bibitem{Yu} X. Yu, \emph{Unbounded runs of non-log-concavity in independence
polynomials of trees}, Zenodo preprint, version 1.1.0 (3 September 2026).
\url{https://zenodo.org/records/22263331}. Main manuscript consulted.
\bibitem{Fang} E. X. Fang, J. Lu, E. Nevo, Y. Yao and H. Zheng,
\emph{Unimodality of Independence Polynomials for Sufficiently Large
Forests}, arXiv:2609.20961v1 (17 September 2026), Section 9.
\url{https://arxiv.org/abs/2609.20961}.
\end{thebibliography}
\end{document}
